\documentclass{aa}

\usepackage{txfonts}

\begin{document}

\title{Efecto Casimir dinámico en cavidad con espejo oscilante}

\subtitle{Informe físico/teórico para la simulación interactiva actual}

\author{Equipo de simulación}

\institute{Proyecto de simulación de física con p5.js}

\date{\today}

\abstract
{El efecto Casimir dinámico aparece cuando las condiciones de borde electromagnéticas cambian en el tiempo, por ejemplo al oscilar uno de los espejos de una cavidad.}
{Describir el marco físico mínimo que explica por qué la simulación muestra aparición de luz y aumento de intensidad al incrementar la rapidez del espejo móvil.}
{Se modela una cavidad de dos espejos paralelos, uno fijo y otro oscilante, y se interpreta la emisión como conversión de fluctuaciones del vacío en fotones reales por modulación temporal de modos.}
{El escenario sin aceleración funciona como referencia de baja emisión, mientras que el modo acelerado favorece un crecimiento más notorio de la señal luminosa en la cavidad.}
{El enfoque es pedagógico y no incorpora pérdidas materiales, ruido instrumental ni un tratamiento electrodinámico completo de detección.}

\keywords{efecto Casimir dinámico -- cavidad electromagnética -- espejo oscilante -- creación de fotones}

\maketitle

\section{Introducción}

El caso estático del efecto Casimir describe una fuerza entre fronteras neutras por modificación del espectro de modos del vacío. Sin embargo, la simulación actual representa el régimen dinámico: dos espejos paralelos forman una cavidad y uno de ellos oscila en el tiempo.

Cuando la longitud efectiva de la cavidad se modula, los modos del campo ya no son estacionarios. Esa variación temporal puede amplificar ciertas fluctuaciones y generar fotones reales, observables como luz en la cavidad.

Para una simulación didáctica, este esquema permite conectar una idea cuántica (vacío fluctuante) con un observable macroscópico (intensidad luminosa), sin introducir toda la complejidad experimental.

\section{Modelo físico}

Se considera un modelo mínimo con los siguientes supuestos:
\begin{itemize}
  \item dos espejos paralelos ideales (uno fijo y uno móvil);
  \item oscilación controlada del espejo móvil en torno a una separación media;
  \item cavidad en vacío y descripción efectiva de la intensidad emitida;
  \item ausencia de pérdidas, ruido y no idealidades de detección.
\end{itemize}

Bajo estas hipótesis, la modulación temporal de la frontera móvil modifica las frecuencias de resonancia de la cavidad y habilita procesos efectivos de creación de pares de fotones.

\section{Ecuaciones clave}

Una forma simple de describir la longitud de cavidad es
\begin{equation}
L(t)=L_0+\delta L\,\sin(\Omega t),
\end{equation}
donde $L_0$ es la separación media, $\delta L$ la amplitud y $\Omega$ la frecuencia de oscilación del espejo móvil.

En un tratamiento efectivo de simulación, la tasa de emisión o intensidad relativa de luz crece con la rapidez de modulación de la frontera. Esto se representa de forma cualitativa como
\begin{equation}
I \propto f(v,\Omega,\delta L),
\end{equation}
donde $v$ resume la rapidez del espejo móvil y $f$ es creciente en el régimen acelerado.

La comparación central usada en la simulación es
\begin{equation}
I_{\mathrm{acelerado}} > I_{\mathrm{no\ acelerado}}.
\end{equation}
Esta desigualdad se evalúa manteniendo iguales los parámetros de referencia.

\section{Interpretación para la simulación}

La simulación representa dos espejos y visualiza brillo dentro de la cavidad cuando el espejo móvil oscila. El objetivo didáctico no es una calibración absoluta de potencia, sino observar tendencias robustas:
intensidad casi nula o baja en referencia no acelerada y crecimiento progresivo de la señal al aumentar la rapidez en modo acelerado.

También se usan magnitudes normalizadas (por ejemplo $I/I_{\max}$) para facilitar comparación entre corridas y evitar sobreinterpretar unidades arbitrarias del entorno gráfico.

\section{Límites del modelo}

Este marco teórico es deliberadamente ideal y tiene límites claros:
\begin{itemize}
  \item \textbf{Pérdidas de cavidad:} disipación y acoplos externos reducen emisión observable.
  \item \textbf{No idealidad del espejo móvil:} amplitud/frecuencia reales no siguen oscilación perfecta.
  \item \textbf{Detección experimental:} ruido de fondo y sensibilidad instrumental afectan la medida de luz.
  \item \textbf{Modelo efectivo:} no reemplaza una electrodinámica cuántica completa del sistema abierto.
\end{itemize}

Por lo tanto, la simulación no debe interpretarse como herramienta de predicción experimental precisa, sino como aproximación conceptual controlada.

\section{Conclusión}

El informe queda alineado con la simulación actual de Casimir dinámico: dos espejos, uno oscilante, aparición de luz en cavidad y aumento de intensidad con rapidez de movimiento.
La comparación entre modo acelerado y no acelerado cumple el rol pedagógico central, al mostrar cómo fronteras dependientes del tiempo pueden convertir fluctuaciones del vacío en fotones observables.
Las simplificaciones adoptadas quedan explícitas para evitar extrapolaciones experimentales indebidas.

\end{document}
